\noindent
This section reviews the literature needed to justify the system requirements and evaluation approach. It covers haptic interfaces for robotics, rendering and contact methods, stability under delay, communication timing, and low-cost hardware platforms.

\subsection{Haptic Interfaces for Robotics}

As established in Section~\ref{sec:introduction_background}, haptic interfaces provide force feedback that complements visual information and supports fine motor control in robotic manipulation and training contexts. This section reviews the relevant literature on rendering techniques, stability analysis, and low-cost hardware platforms.

Commercial haptic devices such as the 3D Systems \emph{Geomagic Touch} and \emph{Touch X}, the \emph{Force Dimension} product range, and the Haply \emph{Inverse3 X} are widely used in research, training, and simulation contexts \cite{giriApplicationBasedReviewHaptics2021}. These systems provide kinesthetic force feedback and are typically distributed with proprietary SDK integration workflows, although some also expose lower-level APIs for custom software development.

However, they are generally sold at substantially higher price points than low-cost academic prototypes. For example, the Geomagic Touch and Touch X are publicly listed at approximately US\$3{,}400 and US\$13{,}400 respectively \cite{geomagicTouchPrice,geomagicTouchXPrice}, while other systems are commonly distributed via quotation-based sales channels \cite{forceDimensionSales,haplyInverse3X}. In contrast, the system developed in this work has a total hardware cost of around £180, highlighting the potential for low-cost systems to provide accessible haptic functionality for research and training applications.

\subsection{Haptic Rendering Techniques}

Real-time haptic rendering has been extensively studied, with particular focus on contact-based interaction models. Two widely adopted approaches are virtual coupling and the virtual proxy (or ``god-object'') method introduced by Zilles and Salisbury and later extended by Ruspini \emph{et al.} \cite{zillesConstraintbasedGodobjectMethod1995,ruspiniHapticDisplayComplex1997}. 

In these approaches, the device state is coupled to a constrained proxy within the virtual environment, enabling stable rendering of contact and geometric constraints. Reliable performance typically requires high update rates on the order of 1~kHz, together with accurate tracking of device state and consistent mapping from Cartesian interaction forces to actuator torques \cite{ruspiniHapticDisplayComplex1997}. In robotic implementations, these mappings are commonly derived using the manipulator Jacobian within operational-space control frameworks \cite{khatibUnifiedApproachMotion1987a}.

Haptic systems are commonly classified as impedance-type or admittance-type devices, depending on whether the device measures position and outputs force (impedance) or measures force and outputs motion (admittance). Most grounded haptic interfaces, including the system developed in this work, operate as impedance devices due to the relative ease of position sensing and actuation.

Contact detection must also be computationally simple enough for the haptic servo loop. Signed distance fields (SDFs) provide direct access to penetration depth and surface normal information from the sign and gradient of the distance function \cite{ericsonRealTimeCollisionDetection2004}. Compared with explicit mesh-contact tests or broad-phase bounding-volume methods alone, this representation is well suited to a 1~kHz proxy-based loop because the contact query produces the quantities required for both proxy projection and normal-force calculation.

\subsection{Stability in Haptic Systems}

Stable haptic interaction imposes strict requirements on latency, control architecture, and numerical implementation. Adams and Hannaford demonstrated that virtual coupling can provide stability guarantees by decoupling device dynamics from the virtual environment under passive interaction assumptions \cite{adamsStableHapticInteraction1999}. 

More generally, stability in haptic systems is often analysed using passivity theory. Colgate and Schenkel showed that discrete-time haptic interfaces can become unstable when energy is artificially generated within the feedback loop due to sampling, delay, or numerical approximation errors \cite{colgatePassivityClassSampleddata1997}. Their passivity framework establishes conditions under which a haptic system remains energetically passive, thereby ensuring stable interaction with arbitrary passive users and environments.

These results highlight the sensitivity of haptic systems to delay and discrete-time effects. In practical implementations, communication latency and asynchronous computation introduce additional energy into the system, reducing stability margins and limiting the achievable virtual stiffness. Passivity-based analyses therefore motivate treating loop delay as a direct limit on renderable stiffness: as delay increases, the range of stable virtual coupling parameters becomes more restricted. This is particularly relevant in systems where communication delay is non-negligible relative to the control timestep, motivating the use of control structures that explicitly manage dynamic coupling, damping, and energy injection in order to maintain stable interaction.

\subsection{Communication Timing in Real-Time Haptic Loops}

Haptic rendering places unusually strict requirements on communication timing because the device state and returned force command form part of the same feedback loop. Prior work on sampled-data haptic interaction shows that delay and discretisation reduce stability margins \cite{colgatePassivityClassSampleddata1997,adamsStableHapticInteraction1999}, while USB latency studies show that USB-based real-time control links can introduce delays on the order of milliseconds \cite{ramadossStudyUniversalSerial2008}. Interfaces that are convenient for prototyping, such as USB virtual COM ports using the USB CDC class, can therefore become a limiting factor when packets are buffered by the interface firmware, host driver, or operating system scheduler. Even when the embedded controller transmits samples at a regular rate, non-uniform host-side delivery can increase effective feedback delay if the control loop processes stale queued states.

For this reason, communication design should be considered alongside control design in practical haptic systems. A latest-value, or snapshot-based, transfer pattern is one suitable architectural response: the haptic loop consumes the newest available state rather than draining a backlog of historical samples. This does not remove transport delay, but it can prevent buffered packets from accumulating into additional application-level latency.

\subsection{Low-Cost Haptic Hardware}

Advances in low-cost components, including 3D-printed structures, microcontrollers, and magnetic encoders, have enabled the development of accessible haptic devices using widely available hardware. This has led to increasing interest in open and modular platforms for research and education \cite{giriApplicationBasedReviewHaptics2021}.

The \emph{Hapkit} platform developed by Martinez \emph{et al.} demonstrates that functional haptic devices can be realised at very low cost (approximately US\$50) using 3D-printed components \cite{martinez3DPrintedHaptic2016}. Subsequent work extended this approach to modular multi-degree-of-freedom configurations, including both pantograph and serial mechanisms \cite{martinezOpenSourceModular2017}. 

Similarly, Lavatelli \emph{et al.} presented an open-source low-cost 2-DOF haptic device, demonstrating that multi-degree-of-freedom actuation can be achieved within an academic budget while maintaining useful interaction capability \cite{lavatelliDesignOpenSource2014}.

However, much of this work focuses primarily on hardware design or educational platforms. Comparatively less emphasis is placed on fully integrated systems that combine low-cost hardware with real-time communication, simulation coupling, and extensible software architectures.

\subsection{Summary of Literature}

The literature establishes that haptic feedback improves performance in robotic manipulation tasks and provides well-developed methods for stable force rendering, particularly through virtual coupling and passivity-based control frameworks. Low-cost hardware platforms have also demonstrated that accessible haptic devices are feasible.

However, there remains a specific gap between simplified hardware demonstrations and fully integrated real-time systems in which communication architecture, stability strategy, simulation coupling, and force rendering are evaluated together. This gap directly motivates the co-designed system architecture presented in Section~\ref{sec:implementation}.
